% Published replacement for the original H5 quotient-model input Q.
\subsection{A published identification of the five-contact quotient}
Let $X=X_0(210,1)$ be the canonical maximal-level Shimura curve of
discriminant $210$, and let $W$ be its Atkin--Lehner group.
Gonz\'alez--Rotger \cite[Theorem 4.2 and Table 2, pp.~938--939]{GonzalezRotger2006}
identify the following curves over $\mathbb Q$ with
$X/\langle w_m\rangle$, respectively:
\[
\begin{array}{c|l}
m&y^2\\\hline
30&-43x^4-686x^3-2915x^2-1372x-172\\
42&-43x^4+600x^3-986x^2-600x-43\\
70&-43x^4+256x^3-130x^2-768x-387\\
105&-43x^4+2x^3+505x^2+12x-1548.
\end{array}
\]
Their notation in Section 4, p.~937, is
$X_D=X_0(D,1)$ and $X_D^{(m)}=X_D/\langle\omega_m\rangle$.
Thus these are the required labeled quotients, not merely curves in
the isogeny classes of their Jacobians. The cited result is a theorem
in the 2006 journal article; no later unpublished quotient table is
needed here.

For a quartic with coefficients $(a,b,c,d,e)$ set
$I=12ae-3bd+c^2$ and
$J=72ace+9bcd-27ad^2-27b^2e-2c^3$.
Its Jacobian has $j=6912I^3/(4I^3-J^2)$. Direct rational arithmetic
on the four published models gives, in the displayed order,
\[
\frac{191342053882402567201}{129708022500},\quad
\frac{2179252305146449}{66177562500},\quad
\frac{128031684631201}{9922500},\quad
\frac{1169975873419524361}{108425318400}.
\]
Writing $j_L(t)=256(1-t+t^2)^3/(t^2(1-t)^2)$, these equal
\[
 j_L(1/2401),\quad j_L(32/375),\quad
 j_L(1/225),\quad j_L(5/1029).
\]
The accompanying exact checker also computes the same values from
the independent short-Weierstrass coefficients of the Jacobians
in Gonz\'alez--Rotger's Table 3, p.~939.

Here is why these invariants determine the desired contact line.
The genus formula in \cite[Proposition 5.2, pp.~940--941]{GonzalezRotger2006}
gives $g(X)=1+\varphi(210)/12=5$: its two elliptic terms vanish
because $5$ splits in $\mathbb Q(i)$ and $7$ splits in
$\mathbb Q(\sqrt{-3})$. Moreover,
$W\simeq(\mathbb Z/2\mathbb Z)^4$.
Each displayed genus-one quotient gives eight fixed points of $w_m$
on $X$. A stabilizer of a point of a characteristic-zero curve is
cyclic, so nontrivial stabilizers here have order two and each
nonempty fixed set is a union of eight-element $W$-orbits.
Riemann--Hurwitz for $X\to X/W$ therefore reads
\[
 8=16(2g(X/W)-2)+\sum_{1\ne w\in W}\#\operatorname{Fix}(w).
\]
The four known fixed sets contribute 32. Hence $g(X/W)=0$, the sum
is 40, and there is exactly one further branch orbit. The product
of the inertia generators in this abelian cover is one. The four
known inertia labels have square-class product $210$; consequently
the fifth label is $210$.

All five labels are distinct and their involutions are rational, so
each branch point is rational: Galois must preserve its unique inertia
label. Normalize their coordinates by
$b_{210}=\infty$, $b_{105}=0$, $b_{70}=1$ and write
$b_{42}=\lambda$, $b_{30}=\mu$.
For $E=X/\langle w_m\rangle$, the residual group has order eight
and exponent two. Over $\overline{\mathbb Q}$ its translation subgroup
is precisely $E[2]$, of order four. The quotient by translations is
isomorphic to $E$ via multiplication by two. Its remaining double
cover of $X/W$ branches at the four points other than $b_m$.
It follows that
\[
\begin{split}
j_L(\lambda)&=j_L(1/2401),\qquad
j_L(\mu)=j_L(32/375),\\
j_L(\mu/\lambda)&=j_L(1/225),\qquad
j_L((\mu-1)/(\lambda-1))=j_L(5/1029).
\end{split}
\]
The first two equations restrict each parameter to its six
cross-ratio transforms. Testing their 36 pairs in the last two
equations gives the unique solution
\[
\lambda=2401,\qquad\mu=-343/32.
\]
This finite computation is performed over $\mathbb Q$ in
\texttt{check\_published\_quotients.py}. Thus the contact configuration
is exactly $\{\infty,0,1,2401,-343/32\}$ from a published modular
identification. This closes input Q only; the arithmetic height
and integral-model arguments are separate obligations.
